\documentclass[11pt]{article}
\usepackage{preamble}
\setlength{\parskip}{4pt}

\begin{document}

\daybanner{AP Statistics \textperiodcentered\ Unit 2 \textperiodcentered\ Topic 2.11 \textperiodcentered\ B: 2026-11-10 \textperiodcentered\ E: 2026-12-04 \textperiodcentered\ TEACHER KEY}{Same Score, Different Standing}

\sectionbanner{CED TETHER}

{\small
\begin{itemize}[leftmargin=1.5em,itemsep=2pt,topsep=2pt]
  \item SKILLS: Using Probability and Simulation
  \item \textbf{Skill 3.A} --- Determine relative frequencies, proportions, or probabilities using simulation or calculations.
  \item \textbf{Skill 3.C} --- Describe probability distributions.
  \item ENDURING UNDERSTANDING: VAR-6 The normal distribution may be used to model variation.
  \item VAR-6.A Calculate the probability that a particular value lies in a given interval of a normal distribution. [Skill 3.A]
  \item VAR-6.A.1 A continuous random variable is a variable that can take on any value within a specified domain. Every interval within the domain has a probability associated with it. VAR-6.A.2 A continuous random variable with a normal distribution is commonly used to describe populations. The distribution of a normal random variable can be described by a normal, or "bell-shaped," curve. VAR-6.A.3 The area under a normal curve over a given interval represents the probability that a particular value lies in that interval.
  \item VAR-6.B Determine the interval associated with a given area in a normal distribution. [Skill 3.A]
  \item VAR-6.B.1 The boundaries of an interval associated with a given area in a normal distribution can be determined using z-scores or technology.
  \item VAR-6.C Determine the appropriateness of using the normal distribution to approximate probabilities for unknown distributions. [Skill 3.C]
  \item VAR-6.C.1 Normal distributions are symmetrical and "bell-shaped." As a result, normal distributions can be used to approximate distributions with similar characteristics.
\end{itemize}
}
\textbf{Essential question:} How can the same raw score have different relative standings in two normal distributions with the same mean?\par
\textbf{DOK-3 skill:} 4.B \textperiodcentered\ \textbf{FRQ pattern:} {\small\texttt{same-\allowbreak{}score-\allowbreak{}two-\allowbreak{}models-\allowbreak{}compare-\allowbreak{}with-\allowbreak{}z}}\par
\textbf{DOK rationale:} Part (c) coordinates two standardized scores and normal areas to adjudicate a same-percentile claim, explain the role of spread, and define what a relative comparison can support.\par

\frameworkphaseheader{Do Now (first take) $\rightarrow$ Explore (video + follow-along) $\rightarrow$ Exit (finish + turn in)}{1 $\rightarrow$ 3}{45}{%
\begin{itemize}[leftmargin=1.5em,itemsep=2pt,topsep=2pt]
  \item Board slide up at the bell. Ask for a claim before students calculate either percentile.
  \item Begin the combined 5.1--5.2 video follow-along at minute 5.
  \item Return to the board slide at minute 33. Circulate and require both z-scores before a comparison.
  \item Collect at the door. Score part (c) only, E/P/I.
\end{itemize}
}{%
\begin{itemize}[leftmargin=1.5em,itemsep=2pt,topsep=2pt]
  \item Choose whether the percentile is the same and cite mean or standard deviation.
  \item Complete the follow-along about sampling variability, normal probabilities, and normal-model appropriateness.
  \item Finish (a)--(c), revise the first take in the exit line, and turn in the sheet.
\end{itemize}
}{%
\begin{itemize}[leftmargin=1.5em,itemsep=2pt,topsep=2pt]
  \item How many School A standard deviations above 80 is a 92?
  \item Why does School B's wider curve place more area above the same score?
  \item Which area gives a percentile: left of 92 or right of 92?
  \item What exactly do you mean by `better' when the raw score is identical?
\end{itemize}
}{Solo teacher. Circulate ~5 pairs. Require students to translate raw-point distance into standard-deviation units and define better as relative standing before choosing a school.}

\begin{callout}{\IconWarn\ WATCH FOR}{calloutred}
\begin{itemize}[leftmargin=1.5em,itemsep=2pt,topsep=2pt]
  \item Using the shared mean as evidence that all corresponding percentiles match.
  \item Reporting the upper-tail probability rather than the percentile below 92.
  \item Reversing the comparison and assigning the higher percentile to the wider distribution.
\end{itemize}
\end{callout}

\sectionbanner{SCORING PART (c) --- E / P / I}

\textbf{Expected elements}
\begin{itemize}[leftmargin=1.5em,itemsep=2pt,topsep=2pt]
  \item \textbf{rejects-same-percentile} (required) --- Rejects the claim that a score of 92 has the same percentile in the two models
  \item \textbf{uses-two-standardizations} (required) --- Cites z approximately 1.33 and percentile 90.9 percent for School A, and z 0.80 and percentile 78.8 percent for School B
  \item \textbf{explains-spread-mechanism} (required) --- Explains that the same 12-point distance represents different numbers of standard deviations because the standard deviations differ
  \item \textbf{makes-relative-choice} (required) --- Chooses School A for the higher relative standing and limits better to the stated percentile meaning
  \item \textbf{states-reader-consequence} (required) --- Explains that the counselor's claim hides the effect of variability on standardized position
\end{itemize}

{\small\renewcommand{\arraystretch}{1.25}
\noindent\begin{tabular}{|p{0.5in}|p{5.9in}|}
\hline
\textbf{E} & Rejects the claim, uses both correct z-scores and percentiles, links the difference to standard deviation, identifies School A as the higher relative standing, and explains the misleading implication. \\ \hline
\textbf{P} & Reaches the correct school and gives both percentiles, but incompletely explains the standard-deviation mechanism, the meaning of better, or the consequence of the claim. \\ \hline
\textbf{I} & Says equal means force equal percentiles, compares only raw distances, or reverses which school gives 92 the higher percentile. \\ \hline
\end{tabular}}

\textbf{Common mistakes}
\begin{itemize}[leftmargin=1.5em,itemsep=2pt,topsep=2pt]
  \item Using 80 as the standard deviation because it is shared by both models
  \item Computing $z=(80-92)/\sigma$ and reporting lower-tail areas without correcting direction
  \item Saying the wider distribution makes every above-mean score more exceptional
  \item Calling one raw score objectively higher even though the numerical score is 92 at both schools
\end{itemize}

\clearpage
\focusbanner{TODAY'S PROBLEM --- annotated}

Suppose exam scores at School A and School B are each approximately normal. Both distributions have mean $80$ points, but School A has standard deviation $9$ points and School B has standard deviation $15$ points.\par\smallskip A counselor says, ``A score of 92 is at about the 90th percentile at both schools because both schools have the same mean.'' A data coach says, ``The different spreads could put 92 at different percentiles.''\par
\begin{center}
\begin{tikzpicture}
\begin{axis}[
  width=0.88\linewidth,
  height=1.75in,
  xmin=35, xmax=125,
  ymin=0, ymax=0.050,
  axis x line=bottom,
  axis y line=left,
  ytick=\empty,
  xtick={50,65,80,92,95,110},
  xlabel={Exam score (points)},
  tick label style={font=\small},
  label style={font=\small},
  legend style={font=\small,draw=none,fill=none,at={(0.98,0.98)},anchor=north east}
]
\addplot[dayblue,very thick,domain=35:125,samples=181]
  {1/(9*sqrt(2*pi))*exp(-0.5*((x-80)/9)^2)};
\addplot[warmred,very thick,dashed,domain=35:125,samples=181]
  {1/(15*sqrt(2*pi))*exp(-0.5*((x-80)/15)^2)};
\addplot[black,densely dotted] coordinates {(92,0) (92,0.049)};
\legend{{School A: $\mu=80,\ \sigma=9$},{School B: $\mu=80,\ \sigma=15$}}
\end{axis}
\end{tikzpicture}
\end{center}
\begin{firsttakebox}{FIRST TAKE (5 min)}
Is a score of 92 at the same percentile in both schools? Choose a claim and name one parameter behind your choice.\par
\answer{Not graded. Equal means should make the counselor's claim tempting; standardization should change the final judgment.}
\end{firsttakebox}

\textit{Rules callout ``STANDARDIZE BEFORE COMPARING (keep on the page)'' is printed on the student sheet.}\par\medskip

\dokbadge{1}~\textbf{(a)} Identify which school has the more variable score distribution and state how the corresponding normal curve shows that difference.\par
\answer{School B is more variable because $\sigma_B=15$ points is larger than $\sigma_A=9$ points. Its normal curve is wider and lower; School A's is narrower and taller around the common mean of 80.}

\dokbadge{2}~\textbf{(b)} For a score of 92, calculate the $z$-score and percentile at each school. Show the standardization and report each percentile to the nearest tenth of a percent.\par
\answer{School A: $z=(92-80)/9=1.333\ldots$, so $P(X_A\leq92)=\Phi(1.333\ldots)\approx0.9088$ and 92 is at about the $90.9$th percentile. School B: $z=(92-80)/15=0.80$, so $P(X_B\leq92)=\Phi(0.80)\approx0.7881$ and 92 is at about the $78.8$th percentile.}

\dokbadge{3}~\textbf{(c)} Decide whether the counselor's same-percentile claim is warranted. Cite both $z$-scores and percentiles, explain why equal means do not give equal relative standing when standard deviations differ, and state at which school a 92 is `better' if better means a higher standing relative to that school's scores. Explain what the counselor's claim would mislead a reader into believing.\par
\answer{The counselor's claim is not warranted. Although 92 is 12 points above the common mean, that gap is $1.33$ School A standard deviations but only $0.80$ School B standard deviations. The corresponding percentiles are about $90.9$ and $78.8$, not the same. If `better' means higher relative standing, the 92 is better at School A. Saying the percentiles are equal would mislead a reader into treating distance in raw points as distance in standard-deviation units and into overlooking School B's greater variability.}

\begin{summaryexitbox}{EXIT}
One thing the video changed about my first take: \hrulefill
\end{summaryexitbox}

\end{document}
